%=====================================================
\section{Sets}
%=====================================================

%-----------------------------------------------------
\begin{frame}{What is a Set?}

\begin{block}{Definition}

\vspace{0.1cm}

A \textbf{set} is a well-defined, unordered collection of distinct objects.
\end{block}

\vspace{0.3cm}

\textbf{Mathematical Representation}

\[
A=\{x\mid P(x)\}
\]

where \(P(x)\) is a predicate that determines whether
\(x\) belongs to the set.

\vspace{0.3cm}

\textbf{Example}
\[
F=\{\text{Age},\text{Income},\text{Experience},\text{Education}\}
\]

represents the feature set used to train a machine learning model.

\end{frame}

%-----------------------------------------------------

\begin{frame}{Representations of Sets}

A set can be represented in two standard forms.

\vspace{0.3cm}

\begin{columns}[T,onlytextwidth]

\column{0.49\textwidth}

\begin{block}{Roster Form}

Lists all elements of the set.

\begin{flushleft}
$\{1,2,3,4,5\}$
\end{flushleft}

\end{block}

\column{0.49\textwidth}

\begin{block}{Set Builder Form}

Defines the set using a common property.

\begin{flushleft}
$\{x\in\mathbb N \mid 1\le x\le5\}$
\end{flushleft}

\end{block}

\end{columns}

\vspace{0.5cm}

\[
\boxed{
\{1,2,3,4,5\}
=
\{x\in\mathbb N \mid 1\le x\le5\}
}
\]

\end{frame}